% -*- coding: utf-8 -*-


\begin{zhaiyao}
\begin{spacing}{1.5}
{

深度神经网络在图像分类等任务中取得了优异性能，但其对对抗样本的高度敏感引发了安全性担忧。现有对抗性训练方法主要针对$L_p$范数约束下的梯度攻击进行防御，对遮蔽类攻击的防御研究相对匮乏。本文围绕基于特征归因的遮蔽攻击及其对抗性训练方法展开研究。首先，实现了两种特征归因方法，即原始梯度显著性图与积分梯度归因方法，分别设计了固定遮蔽攻击和自适应遮蔽攻击两种攻击策略，并对比了两种归因方法在攻击效果与计算效率上的差异。在MNIST数据集上基于LeNet-5模型进行了系统实验。结果表明，积分梯度遮蔽攻击在攻击强度上优于原始梯度遮蔽攻击，但计算代价更高；基于遮蔽攻击的对抗性训练可有效提升模型对遮蔽攻击的防御能力，但可能降低模型对梯度攻击的防御能力；混合对抗性训练策略实现了梯度攻击防御与遮蔽攻击防御的较好平衡。

}

\end{spacing}
\end{zhaiyao}




\begin{guanjianci}
对抗样本，对抗性训练，显著性图，梯度积分，遮蔽攻击
\end{guanjianci}



\begin{abstract}
\begin{spacing}{1.5}
Deep neural networks have achieved remarkable performance in tasks such as image classification, yet their vulnerability to adversarial examples raises serious security concerns. Most existing adversarial training methods are tailored to defending against gradient-based attacks under an $L_p$-norm constraint, while defenses against occlusion attacks remain relatively underexplored. This thesis investigates saliency map-based occlusion attacks and corresponding adversarial training methods. Specifically, we implement two feature attribution approaches, gradient saliency maps and Integrated Gradients, and design two attack variants: a fixed occlusion attack and an adaptive occlusion attack. We then systematically compare these two attribution methods in terms of attack effectiveness and computational cost, and conduct experiments on the MNIST dataset using LeNet-5. The results indicate that Integrated Gradients-based occlusion attacks achieve stronger attack strength than saliency-based occlusion attacks but incur substantially higher computation. Adversarial training with occlusion attacks can significantly improve robustness to occlusion attacks, but may reduce robustness to gradient-based attacks. A mixed adversarial training strategy provides a better balance between robustness to gradient attacks and to occlusion attacks.

\end{spacing}
\end{abstract}


\begin{keywords}
Adversarial Examples, Adversarial Training, Saliency Maps, Integrated Gradients, Occlusion Attacks
\end{keywords}
